El objetivo principal del proyecto es el desarrollo de una aplicación web que permita centralizar y digitalizar los principales procesos asociados a la actividad comercial del distribuidor, facilitando la gestión de pedidos, el seguimiento de la actividad y la relación con los clientes profesionales.

De manera más específica, los objetivos del sistema son los siguientes:

\begin{itemize}
\item Facilitar la gestión de pedidos entre el distribuidor y las peluquerías, permitiendo consultar el catálogo de productos y realizar pedidos de forma sencilla.

\item Mejorar el control y seguimiento de la actividad comercial del distribuidor, incluyendo aspectos como la gestión de rutas, el registro de cobros o el seguimiento de pedidos.

\item Reducir la dependencia de procesos manuales y herramientas dispersas, centralizando la información relevante del negocio en una única plataforma.

\item Proporcionar herramientas digitales útiles para las peluquerías, como el acceso al catálogo de productos o funcionalidades de apoyo a la gestión de su actividad profesional.

\item Mejorar la comunicación y relación entre el distribuidor y sus clientes, favoreciendo la fidelización y el uso continuado del sistema.

\end{itemize}

En conjunto, se espera que la implantación del sistema contribuya a mejorar la eficiencia operativa del distribuidor, reducir el tiempo dedicado a tareas administrativas y facilitar la gestión de la actividad comercial. Asimismo, el desarrollo realizado sienta las bases técnicas necesarias para la incorporación progresiva del resto de funcionalidades previstas en el sistema.
